\lecture{09}{Sep 12 10:15}{Equivalent Systems}
\section{Equivalent Systems}
\begin{definition}[Equivalent]\label{def:equivalent}
    Two linear systems are called \textbf{equivalent} if their solution sets are identitical.
\end{definition}
\begin{definition}[Elementary Change]\label{def:elementary-change}
    An \textbf{elementary change} to a linear system is one of the follow changes:
    \begin{enumerate}
        \item Swapping the position of two equations
        \item Multiplying every coefficient in one equation by a non-zero scalar
        \item Adding the same multiple of the coefficients of one equation to the respective coefficients of another equation
    \end{enumerate}
\end{definition}
\begin{proposition}\label{prop:elementary-change}
    Performing an elementary change on a linear system does not change the system's solution set.
\end{proposition}


\begin{proof} Type (i)
Let $S$ be the solution set of a linear system and let $S'$ be the solution set of the system following the elementary change $R_i \leftrightarrow R_j$.

The system includes the same equations before and after the change, so clearly $S = S'$.
\end{proof}
\begin{proof} Type (ii)

Let $S$ be the solution set of a linear system and let $S'$ be the solution set of the system following the elementary change $R_i \to tR_i$, where $t \neq 0 \in F$.

All equations except the $i$-th remain unchanged. For the $i$-th equation:

Suppose the $n$-tuple $c = (c_1, c_2, \ldots, c_n) \in S$, then
\[a_{i1}c_1 + a_{i2}c_2 + \cdots + a_{in}c_n = b_i\]

Therefore,
\[ta_{i1}c_1 + ta_{i2}c_2 + \cdots + ta_{in}c_n = t(a_{i1}c_1 + a_{i2}c_2 + \cdots + a_{in}c_n) = tb_i\]

Therefore $c$ is also a solution of the $i$-th equation in the changed system, and $c \in S'$.

Conversely, suppose $c \in S'$, then
\[ta_{i1}c_1 + ta_{i2}c_2 + \cdots + ta_{in}c_n = tb_i\]

Therefore, because $t \neq 0$,
\begin{align*}
a_{i1}c_1 + \cdots + a_{in}c_n &= t^{-1}(ta_{i1}c_1 + ta_{i2}c_2 + \cdots + ta_{in}c_n)\\
&= t^{-1}(tb_i)= b_i
\end{align*}

Therefore, $c$ is also a solution of the $i$-th equation in the original system, and $c \in S$.
Since $S \subseteq S'$ and $S' \subseteq S$, therefore $S = S'$.
\end{proof}
\begin{proof} Type (iii)

Let $S$ be the solution set of a linear system and let $S'$ be the solution set of the system following the elementary change $R_i \to R_i + tR_j$ for $t \in F$.

All equations except the $i$-th remain unchanged. For the $i$-th equation:

Suppose the $n$-tuple $c = (c_1, \ldots, c_n) \in S$ is a solution of the original system, then
\[a_{i1}c_1 + a_{i2}c_2 + \cdots + a_{in}c_n = b_i\]
and
\[a_{j1}c_1 + a_{j2}c_2 + \cdots + a_{jn}c_n = b_j\]

Therefore,
\begin{align*}
&(ta_{j1} + a_{i1})c_1 + \cdots + (ta_{jn} + a_{in})c_n\\
&= ta_{j1}c_1 + \cdots + ta_{jn}c_n + a_{i1}c_1 + \cdots + a_{in}c_n\\
&= t(a_{j1}c_1 + \cdots + a_{jn}c_n) + (a_{i1}c_1 + \cdots + a_{in}c_n)\\
&= tb_j + b_i
\end{align*}

This shows that $c$ is also a solution of the $i$-th equation in the changed system and therefore $c \in S'$ and $S \subseteq S'$.

Conversely, suppose $c \in S'$ is a solution of the changed system, then
\[(a_{i1} + ta_{j1})c_1 + (a_{i2} + ta_{j2})c_2 + \cdots + (a_{in} + ta_{jn})c_n = b_i + tb_j\]
and
\[a_{j1}c_1 + a_{j2}c_2 + \cdots + a_{jn}c_n = b_j\]

Therefore,
\begin{align*}
&a_{i1}c_1 + \cdots + a_{in}c_n\\
&= (a_{i1}c_1 + \cdots + a_{in}c_n) + 0\\
&= (a_{i1}c_1 + \cdots + a_{in}c_n) - ((a_{i1} + ta_{j1})c_1 + \cdots + (a_{in} + ta_{jn})c_n) + (b_i + tb_j)\\
&= (a_{i1}c_1 + \cdots + a_{in}c_n) - (a_{i1}c_1 + \cdots + a_{in}c_n) + b_i - t(a_{j1}c_1 + \cdots + a_{jn}c_n) + tb_j\\
&= 0 + b_i - tb_j + tb_j\\
&= b_i
\end{align*}

This shows that $c$ is also a solution of the $i$-th equation in the original system and therefore $c \in S$ and $S' \subseteq S$.

Therefore, since $S' \subseteq S$ and $S \subseteq S'$, we have $S = S'$.
\end{proof}



\begin{theorem}\label{thm:equivalent-systems}
    A \textbf{finite sequence of elementary changes} on a linear system leads to an \textbf{equivalent} system.
\end{theorem}
\begin{proof}
    Let $S$ be the solution set of a linear system and $S'$ be the solution set of the system after performing a finite sequence of elementary changes.

    By Proposition \ref{prop:elementary-change}, any single elementary change leads to an equivalent system.

    Equivalency is transitive, because equality of sets is transitive.
\end{proof}